
\section{Test Generation Pipeline}
We can use a formal leakage predicate for a defense as input to a generation pipeline to produce a set of litmus tests that induce exactly the $\mu$arch leakage scenarios the defense implementation ostensibly prevents. We will then run this set on the defense hardware implementation \anna{hardware under test??}, to ensure that it blocks all such leaks. We implement a pipeline that produces a set complete across high-level types of leakage with additional random fanout to test implementations robustly while avoiding test set explosion. We  run small experiments and analyze examples to ground our selection of variables to focus on during test generation. 
\anna{delete commented i think:}
% We designate the design dimensions we identify as either \textbf{1. comprehensively enumerated, 2. semi-enumerated, or 3. fixed}, where (2) and (3) are heuristically determined, and expose all dimensions of the pipeline to tuning. 


Fundamentally, our generation pipeline is comprehensive across types of leakage backbones, that is, it will generate all possible leakage backbones as blocked by the defense, but is heuristic-based in the concretization phase, e.g. it does not test every combination of eligible x86 .asm concretization for each instruction. 

The pipeline takes in a defense leakage predicate,  uses it to constrain an execution model and generate violating execution graphs (\ref{stage1}), then concretizes those execution graphs to program-annotation pairs
\urgentbr
\urgentnote{\textbf{STAGE NAMES (Fig.~\ref{fig:pipeline} now says Stage~1a, 1b, 2, 3; renamed 2026-11-11).} This sentence is the one place the text walks through the stages, so name all four here. It also has an undefined reference (\texttt{\textbackslash{}ref\{stage1\}}, no such label; it prints ??) and is unfinished. \textbf{PROPOSED TEXT} (replaces the sentence just before this note): The pipeline takes a defense's leakage predicate and uses it to constrain an execution model in Alloy: Stage~1a enumerates the minimal violating execution graphs, and Stage~1b grows each one into larger graphs under loosened minimality (Fig.~\ref{fig:pipeline}). Stage~2 concretizes each graph into program-annotation pairs, and Stage~3 runs them on the defense's gem5 build and checks for transmission.}
\notebr


% \input{tikz/pipeline_full}
\anna{define leakage backbones}


% Once we have a formal model of a defense's security guarantee, we can use this specification as input to a test generation pipeline to build a set of litmus tests \textit{that are tuned to exhibit microarchitectural leakage} according to those guarantees.
\subsection{Abstract test generation}

Since side-channel leakage is a $\mu$arch behavior, to be comprehensive across the leakage space we enumerate tests first as abstract execution graphs, with abstract instructions and $\mu$arch state annotations that represent a time snapshot of a processor. We aim to enumerate \textit{all possible fundamental forms of Spectre leakage} that are blocked by the target defense, and then concretize these abstract \textit{leakage backbones} 
into (1) assembly tests and (2) corresponding speculation and other $\mu$arch directives with which to non-invasively instrument our simulator to elicit leakage. A \emph{leakage backbone} is a minimal violating execution graph: removing or relaxing any of its instructions, operands, state elements or flags makes it non-violating. First, we must have a format for comprehensive enumeration of the forms of Spectre leakage.

\subsubsection{Test Generation (Alloy)}
We define constraints on execution graphs in the Alloy relational modeling language, so that we can use the Alloy Analyzer SAT-based finite model-finder to solve for all satisfying graphs up to a bound~\cite{jackson2002alloy}. As shown in figure \ref{fig:pipeline}, there are 3 types of constraint:
\urgentbr
\urgentnote{\textbf{STAGE NAMES}: this subsubsection describes Stage~1a, so say so. \textbf{PROPOSED TEXT} (edits the sentence just before this note): As shown in Stage~1a of Fig.~\ref{fig:pipeline}, there are three types of constraint: (and, after the list, the Stage~1a note below already proposes the UNSAT sentence).}
\notebr
\begin{enumerate}
    \item \textbf{Execution well-formedness constraints}: we must enforce that an arbitrary graph of instructions can arise from a program. E.g. the speculative program order (spo) relation must be a total order on instructions. These constraints are defense-agnostic.
\notebr
    \item \textbf{Minimality requirements:} we enforce that our graph is minimal along certain dimensions, that is all components specified in the graph are load-bearing to the execution violating the leakage predicate. This is measured by requiring every \textit{relaxation} of the graph, i.e. removal or weakening of an element, to be non-violating.
\notebr
\auditnote{write}{writing, high}{Decided 2026-11-08: Stage~1a adds instruction minimality (RI) to the models' own RR, RC+RR and RS checks; RO is defined but enforced in neither stage; Stage~1b drops RI again. The text should say which stage checks which relaxation. \newline \textbf{PROPOSED TEXT} (edits the sentence just before this note): In Stage~1a, a graph is minimal if removing any instruction (RI) or any state element (RS), resolving any unresolved instruction (RR), or committing and resolving the first uncommitted instruction (RC) makes it non-violating; Stage~1b keeps RR, RC and RS but not RI. \newline \textit{Source:} STAGE1\_alloy/models/proposed\_twophase/STT\_6\_core\_ri.als:289-302 (gen\_useful\_litmus, RI at :301); STT\_6.als:301-311; URGENT\_TODO.md 'Stage 1 minimality' decision (2026-11-08)}
\notebr
    \item \textbf{Defense leakage predicate:} we construct the defense leakage predicate from the three components in section \ref{formalmodel}, and require our execution to violate the predicate (i.e. be an example of leakage)
\notebr
\end{enumerate}
% The model describes a space of executions partitioned by the leakage predicate, which allows us to 
Alloy synthesizes all execution graphs up to a size bound that satisfy these constraints by iteratively model-finding and blocking discovered graphs until the problem is UNSAT.
\notebr
\auditnote{write}{writing, high}{Partly out of date. Stage~1a now enumerates up to six instructions (`but 6') by model-finding and blocking until UNSAT, and reaches UNSAT for STT and ReCon and their store-bypass variants (STT\_6: 129 cores; STT\_6\_stl: 306; every `rest' coverage partition UNSAT). SPT's Stage~1a does not converge: its four-instruction partition hits the 200,000 cap. Stage~1b enumerates a capped number of graphs per Stage-1a graph (about 300,000 per family). So only the Stage-1a cores are exhaustive, and only for STT and ReCon. The same claim is in the intro (`complete in including all leakage backbones') and in the knobs subsection (`coverage is complete up to the instruction bound'). \newline \textbf{PROPOSED TEXT} (edits the sentence just before this note): Stage~1a enumerates every minimal violating graph of up to six instructions by repeatedly model-finding and blocking each graph found until the problem is UNSAT (for SPT, up to a cap); Stage~1b then enumerates up to a fixed number of six-instruction graphs around each one. \newline \textit{Source:} wiki/state/EVAL\_MANIFEST.md `Stage 1' and `Stage 1.5' sections; STAGE1\_alloy/models/proposed\_twophase/STT\_6\_core\_ri.als:308-310}
\notebr
\subsubsection{Relaxations}
We define a minimality criterion on execution graphs in order to limit the problem solutions so that we can synthesize all \textit{leakage backbones}, or fundamental forms of leakage. Then we can later incorporate complexity by concretizing one backbone into a program in multiple ways, with the assurance that we have at least some representative program for all backbones. This also allows us to group executions that exhibit the same form of leakage, and restrict initial synthesis to the smallest example of each class, to keep our test set compact.
 % and make claims about fundamental bounds on leakage backbone size and quntity
% . We want to group executions that exhibit the same form of leakage, to allow us to reason about the simplest form of this leakage and build out in a principled manner.
% For this reason we follow MCM work~\cite{lustig} to define a minimality criterion restricts synthesis to the smallest example of each fundamental leakage backbone, keeping the initial test set compact while covering every leakage type. We synthesize all \textit{minimal, violating} executions up to a specified instruction bound. 
Our minimality criteria are established by a set of \textbf{relaxations}, which are perturbations on our execution instance, under which the instance should transition to a non-violating state. This ensures that all components of a generated execution are necessary for leakage, allowing us to make claims about the fundamental bounds on leakage backbone size and quantity.
\urgentbr
\urgentnote{\textbf{TODO (your note from the old draft, sections/oldtechnicaldetails.tex:113): add a natural-language intuition proof that every leakage backbone fits within a bounded length}, for STT first, using the running example. What the runs already show: Stage~1a reaches UNSAT at six instructions for STT and ReCon, so every minimal violating graph of up to six instructions is found. They do not show that no backbone needs seven or more instructions; that is what the proof would add. A cheaper empirical check is to rerun Stage~1a with a bound of seven and see whether any new minimal graph appears. Until then, ``claims about the fundamental bounds on leakage backbone size'' should say ``up to the instruction bound''; for SPT, whose Stage~1a stops at its cap, even the up-to-six claim does not hold yet. \textit{Source:} wiki/state/EVAL\_MANIFEST.md `Stage 1' section; the Stage~1a note in the Test Generation subsection above.}
\notebr

Figure~\ref{fig:relaxation} shows examples of tests that get minimized to the same test.
\input{tikz/relaxation}% the one relaxation figure (your four panes, redrawn), 2026-11-03
\notebr
\auditnote{fig}{figure, low}{\textbf{About Fig.~\ref{fig:relaxation}:} Pane (c) puts 'unresolved' on a Load, which breaks the rule that resolution is shown only on branches and stores. The reader is never told that an unflagged load (pane d) counts as resolved. The caption also leaves the dashed rf edge in (b) unexplained. \emph{Proposed:} Add: 'Non-branch rows without a flag word are resolved and uncommitted; dashed = a removed edge.' \emph{Source:} Fig. 4 p.5 render; VISUAL\_LANGUAGE.md CURRENT RULES.}
\notebr
\auditnote{fig}{figure, high}{\textbf{About Fig.~\ref{fig:relaxation}:} With RI-only Stage~1a (decided 2026-11-08) the caption's claim holds for Stage~1a: RI removes (b) and RR removes (c). (a) is removed too, but by RS (r1 appears only in the unneeded operand, and removing a state drops every operand on it), because RO is enforced nowhere. Stage~1b re-admits graphs like (b) by design; in the flat test sets behind the June numbers, (b) was never removed. This comes from reading the model, not from an Alloy run. \newline \textbf{PROPOSED TEXT}: relabel (a)'s arrow RS, or keep RO and add `RO is defined but not enforced; RS removes (a)'; end the caption with `\ldots{} so Stage~1a does not synthesize them.' \newline \textit{Source:} STAGE1\_alloy/models/STT\_6.als:214 (o\_p drops operands on a removed state), :301-311; proposed\_twophase/STT\_6\_core\_ri.als:301}
\notebr
We define the following relaxations:
\begin{enumerate}
    \item \textbf{RI}: remove instruction
    \item \textbf{RO}: remove operand
    \item \textbf{RS}: remove state
    \item \textbf{RC}: shift commit: flip the first uncommitted instruction to committed
\notebr
\auditnote{fact}{fact}{The code applies RC to the first uncommitted instruction, the one right after the commit boundary. 'Last-uncommitted' reads as the youngest instruction. RC is also never used alone: it is tested together with RR on that same instruction. When nothing is committed, first\_uncommitted falls back to the first instruction. \newline \textbf{PROPOSED TEXT} (edits the sentence just before this note): \textbackslash{}textbf\{RC\}: shift commit: mark the first uncommitted instruction (the commit boundary) committed and resolved \newline \textit{Source:} STT\_6.als:306-308 (secure[RC-\textgreater{}first\_uncommitted + RR-\textgreater{}first\_uncommitted]), :139 (fallback to first\_instr)}
\notebr
    \item \textbf{RR}: mark an unresolved instruction as resolved
\end{enumerate}

\subsubsection{Overlapping interleaving}
Once we have defined minimality criteria for an execution, we are then able to utilize all our relaxations to make claims about fundamental bounds on leakage backbones, as discussed in the results. However, we recognize a need for more complex tests, so that we can stimulate exercise of protection set tracking implementation, as well as other implementation mechanisms. For instance, a correct implementation of STT blocks the leak in Figure~\ref{fig:ddi-rf}, while the test in Figure~\ref{fig:interleave} adds a chain of instructions after the committed prefix that merges a second speculative source into an ALU op; this shape exposes the ReCon \texttt{getOldestTaint} bug (Sec.~\ref{sec:bug-otb}), where the implementation lets it leak.
\input{tikz/interleave_v4_8pt}
\notebr
\urgentbr
\urgentnote{\textbf{STAGE NAMES}: this figure's caption (in tikz/interleave\_v4\_8pt.tex) says ``no Stage~1 instance contains''; change it to ``no Stage~1a instance contains''. The claim holds for Stage~1a, whose graphs are minimal; Stage~1b drops instruction minimality, so it is not promised there.}
\notebr
\auditnote{fig}{figure}{\textbf{About the caption of Fig.~\ref{fig:interleave}:} With RI-only Stage~1a, `no Stage~1 instance contains' is now true by construction: the second source is redundant for the leakage predicate, so no minimal graph has it. The mechanism is what is out of date. Chain interleaving is off in every evaluation config, put an ALU on a load transmitter's address path in 0 of 1,200 variants (that address is always an Alloy-locked register, and chains write only free registers), and cannot hook onto a Stage-1a core. SimSpect reproduces the getOldestTaint bug with the OTB merge of fanout v2 instead: a second load under an earlier, resolved branch, xor-merged into the transmitter's address (flat test set, ReCon-STT build: 100 of 100 tests hit, 0 of 100 without the merge). The merge is on in the 2026-11-09 two-stage evaluation configs, but it applies to few two-stage programs (it needs a resolved branch before the unresolved one); where it applies, the ReCon-STT row hits the getOldestTaint drop in 11 of 24 OTB-transmitter tests per mode at the 2\% sample. \newline \textbf{PROPOSED TEXT}: redraw the figure and caption around the merge the evaluation uses, or caption this figure as the chain-interleaving capability rather than the route behind the result. \newline \textit{Source:} wiki/design/VOCAB\_TABLES.md (fanout v2 OTB); runconfigs/twostage\_fv2\_20261109/E2\_STT\_6\_fv2.jsonc (interleave off, fanout\_v2.otb on); URGENT\_TODO.md `OTB on two-stage tests'}
\notebr
\auditnote{fig}{figure}{\textbf{About the caption of Fig.~\ref{fig:interleave}:} Branch flags are inconsistent across the graph figures. Fig. 4, 6, 8, 9 and Table II write 'unresolved' on the branch row; Fig. 5 (i2) and Fig. 10 (i2) write nothing; Fig. 2 puts it only in the note column. Captions of Fig. 5, 6 and 10 also use subscripts (\$f\_0\$, \$r\_0\$, \$r\_1\$) where the figures print f0, r0.  \textbf{PROPOSED TEXT} (edits the sentence just before this note): Write the italic 'unresolved' flag on every unresolved branch (Fig. 2 i0, Fig. 5 i2, Fig. 10 i2), and use r0/f0 without subscripts in the captions.  \textit{Source:} Renders p.4-p.11; VISUAL\_LANGUAGE.md CURRENT RULES ('Resolution status ... shown only on branches and stores'; 'write r0, not \$r\_0\$').}
\notebr
\auditnote{fig}{figure, high}{\textbf{About the caption of Fig.~\ref{fig:interleave}:} The figure marks i0 and i1 as committed, so c0 (placed after them and before i2) is not speculative at all under STT's contract. As drawn, the merge has only one tainted source, and the caption's getOldestTaint mechanism cannot occur. The OTB gadget needs the injected load under a branch that is resolved early but still uncommitted.  \textbf{PROPOSED TEXT} (edits the sentence just before this note): Mark i0/i1 'resolved' instead of 'committed', or explain as Fig. 6 does that i0 is committed only at the observation time and resolves at a swept time.  \textit{Source:} Fig. 5 p.6 render (i0, i1 'committed'); OTB\_GENERATION\_GAP.md 'What OTB needs' (ld\_inj under br\_r, a resolved branch that resolves early); STT\_6.als speculation\_contract\_p.}
\notebr
Therefore, we interleave instructions in two rounds. The first round happens in Alloy, by removing some of the relaxations: the models we evaluate enforce only RR, RC and RS, and leave RI and RO defined but unchecked. The second round happens during concretization (Sec.~\ref{sec:conc-interleave}) and builds merges like the one in Figure~\ref{fig:interleave}.
\urgentbr
\urgentnote{Out of date since the two-stage decision (2026-11-08): the evaluated Stage~1a also enforces RI, and Stage~1b is the round that drops it; RO is enforced in neither. \textbf{PROPOSED TEXT} (edits the sentence just before this note): The first round happens in Alloy: Stage~1a enumerates minimal graphs (RI, RR, RC and RS), and Stage~1b grows each one with RI turned off, so it may add instructions that share the backbone's state (RO is defined but not enforced in either stage). \textit{Source:} STAGE1\_alloy/models/proposed\_twophase/STT\_6\_core\_ri.als:289-302; wiki/state/EVAL\_MANIFEST.md `Stage 1.5'.}
\notebr
The first round allows us to emit tests that include instructions not on the leakage backbone, but to minimize the load on the Alloy Analyzer, we still require all instructions to share state that is necessary to the leakage backbone. We will do further interleaving in the concretization phase for the non-state-sharing instructions. 
% Figure~\ref{fig:statetable} shows the types of instructions and state enumerated in each phase. To enable our first round of interleaving, we turn off the remove instruction relaxation, and the remove operand instruction, and require all operands to use essential state.

\subsection{Concretization Phase}
Given our abstract execution graphs which include $\mu$arch state, we must generate programs and simulator directive combinations to elicit this execution state. Fig.~\ref{fig:pipeline} (Stage~2) lists the concretization passes; we describe them in that order.

\subsubsection{Instruction Choice}
\reviewtag{[TODO: Applied your requests in this subsection; check it]}
First, we concretize our abstract instructions to assembly. We can tune our assembly translation by defining instruction tables which contain all possible candidates for each abstract instruction form. The SimSpect pipeline selects from the candidates that use enough operands to specify the abstract instruction as given in the Alloy execution. Each chosen instruction is emitted as an inline-assembly block with a global PC label
\notebr
\auditnote{fact}{fact, low}{Two qualifications. Branches are not picked at random: pass 2.5 rewrites every non-transmitter branch to br\_cond and the transmitter to br\_bez. Also, for a branch only the condition set-up is inline asm with the PC label. LLVM emits the compare and the jcc, so LLVM's block layout decides which jcc direction is the speculative path. That is why the harness derives the forced direction from the binary. \newline \textbf{PROPOSED TEXT} (edits the sentence just before this note): Each chosen instruction is emitted as an inline-assembly block with a global PC label (for a branch, the label sits on the inline-asm condition set-up and LLVM emits the compare and jump) \newline \textit{Source:} STAGE2\_compilation/parsexml.py:789-795 (branch rewrite), :2426-2427 and :2484 (icmp + br emitted by LLVM); wiki/design/RESOLUTION\_HOOKS.md sec.2}
\notebr
(Table~\ref{tab:instrtable}).
We select uniformly at random among the valid candidates, though we could exhaustively enumerate the cross product of valid concretizations.

\subsubsection{Speculation Primitive Instrumentation}
Table~\ref{tab:instrtable} lists, for each speculation primitive, the program Stage~2 emits (including whether each branch is mispredicted and which way it is taken) and the setting the Stage~3 harness applies.
The concretization pipeline must realize speculation primitives in a way that ensures that a test experiences the correct speculation as seen in the abstract execution. In the model, we introduce resolved and committed flags. At some observation time during execution, the run must match the speculative program order and resolved and committed flags in the modeled instance. An unresolved branch is always mispredicted, in both directions (architecturally taken and mispredicted not taken, or the reverse); its constant condition sends the architectural path to the exit, so the rest of the test runs only speculatively. A resolved branch is generated only as correctly predicted and not taken
\notebr
\auditnote{fact}{fact, low}{There is one exception. The transmitter branch is always emitted as br\_bez: architecturally taken, forced to mispredict, and given a resolve stall of 0, whatever its Alloy resolved flag. Table row 5 says this but the text does not. \newline \textbf{PROPOSED TEXT} (edits the sentence just before this note): A resolved branch is generated only as correctly predicted and not taken; the transmitter branch is the exception and is always mispredicted (Table\~{}\textbackslash{}ref\{tab:instrtable\}, row\~{}5). \newline \textit{Source:} STAGE2\_compilation/parsexml.py:738-776 (xm override); pipeline.py:404-421 (\_inject\_stalls: xmit stall 0)}
\notebr
(Table~\ref{tab:instrtable}, rows 1--2). Every case can induce leakage while the branch is unresolved, even if it is not ultimately mispredicted, but we do not generate correctly predicted unresolved branches at this time for simplicity.
We also sweep the resolution time of resolved branches, by setting a delay from the time they receive their direction and target and the time they transmit ``resolved,'' which we will direct in the gem5 implementation. This allows testing of multiple branches with different resolution orders, which can induce different behaviors (Fig.~\ref{fig:resolutiontiming}); the ReCon bug of Fig.~\ref{fig:otb} is one such case (sweep values: Table~\ref{tab:instrtable}, row~2). For unresolved branches, we direct them to resolve with a fixed 5000-cycle stall, and later enforce that this stall lasts until after our observation time.
\input{tikz/instrtable_v2_8pt}
\notebr
\auditnote{fact}{fact}{\textbf{About the caption of Table~\ref{tab:instrtable}:} No evaluated model can produce an Other\_x transmitter. Otherxs is commented out of the leakage function in STT\_6, Recon\_6 and SPT\_6\_oneLP, and tag\_xm requires the xmit to own a leakage-function operand. So this row, check\_other.py and the divide-fault promotion never ran in the reported campaigns. The same applies to pipeline.tex's 'Faulting / long-latency other transmitter' bullet.  \textbf{PROPOSED TEXT} (edits the sentence just before this note): (caption) ... Other\textbackslash{}\_x rows are supported by the pipeline, but no evaluated defense model lists Other\textbackslash{}\_x as a transmitter.  \textit{Source:} STT\_6.als:327-328, Recon\_6.als:320, SPT\_6\_oneLP.als:344 (leakage\_function); STT\_6.als:270-273 (tag\_xm); STAGE3\_gem5/gem5\_common.py:772-852; OTB\_GENERATION\_GAP.md ('no other\_x transmitters in any STT\_6 testset')}
\notebr
\auditnote{fig}{figure, low}{\textbf{About the caption of Table~\ref{tab:instrtable}:} Table II says 'first uncommitted' where the text and Fig. 7 say 'first non-committed', and it adds an 'uncommitted' flag word that Table I does not define. The 'Target transmitter' row covers only Branch\_x and Other\_x (calling the latter 'Divide'), so readers cannot see that Load/Store transmitters need no special handling, or that Other\_x is a divide or zero-skip multiply.  \textbf{PROPOSED TEXT} (edits the sentence just before this note): Use 'first non-committed'; add to row 5: 'Load, Store: no special code'; say 'Other\_x (divide)'.  \textit{Source:} Table II p.6 render; plans/fig\_d\_instrtable.md Row 5 (kinds: load, store, branch, zsm/divide).}
\notebr
\input{tikz/resolutiontiming}% your figure, redrawn in the new notation, 2026-11-03
\notebr
\auditnote{fig}{figure}{\textbf{About the caption of Fig.~\ref{fig:resolutiontiming}:} The graph violates STT's leakage predicate under both options; what changes is whether a getOldestTaint defense releases i5. Labeling only option 2 'violation' (in red) conflicts with the meaning of 'violation' in Sec. IV and with Fig. 2.  \textbf{PROPOSED TEXT} (edits the sentence just before this note): "(1) i0 resolves before i4 executes: ReCon holds i5" / "(2) ... after: ReCon releases i5 (bug)".  \textit{Source:} Fig. 6 p.7 render; pipeline.tex violation definition; Fig. 6 caption ('the run commits without violating the contract').}
\notebr
% \urgentnote{\textbf{VOCAB}: In Fig.~\ref{fig:resolutiontiming} the graph is a violation under both options; what differs is whether the run leaks. Change the figure's notes ``no violation'' to ``no leak'' and ``violation'' to ``leak'', and in its caption ``the run commits without violating the contract'' to ``the run does not leak''.} 
\urgentbr
\urgentnote{Answer to the comment above (``this makes no sense''): in the paper's terms, \emph{violation} is a property of the abstract graph, and this graph breaks STT's leakage predicate under both timings: $m1$ reaches i5's address through i3 and i4 while i2 is unresolved. The two options differ only in what the gem5 run does. In option~1 the defense holds i5 (no hit); in option~2 the getOldestTaint tracking lets i5 send its access (a hit: the bug). So the figure's notes should name what the run does, not ``no violation''/``violation''. The commented line proposed ``no leak''/``leak''; the hidden figure note above proposes ``ReCon holds i5''/``ReCon releases i5 (bug)''. Either works; pick one and delete the commented line. \textit{Source:} tikz/resolutiontiming.tex:104-105 (the two notes), :109 (caption).}
\notebr
\auditnote{fig}{figure}{\textbf{About the caption of Fig.~\ref{fig:resolutiontiming}:} Fig. 6 and Fig. 10 draw the same gadget (per the source comment) in different ways: Branch\_n vs Branch\_x, r2 \textless{}- r1,r0 vs r0 \textless{}- r1,r0, protected m1 vs m0 with no protection mark. Fig. 10's caption calls i1 'speculative under branch i0' while marking it committed, without Fig. 6's explanation.  \textbf{PROPOSED TEXT} (edits the sentence just before this note): Draw Fig. 10's graph with Fig. 6's kinds, names and protection mark, or say in Fig. 10's caption that i0/i1 are committed only at the observation time (as in Fig. 6).  \textit{Source:} Fig. 6 (p.7) vs Fig. 10 (p.11) renders; resolutiontiming.tex header comment 'the same two-branch merge shape'; RESULTS\_FIGURES.md §7 ('m0 is not gray').}
\notebr
For data-flow speculation, we instrument store bypass~\cite{horn2018ssb} (a store-to-load-forwarding misprediction) only in the false negative direction, where a load mispredicts no alias with an older store to the same address, whose address we hold unresolved, so the load bypasses it and reads the stale value until the memory-order squash. For this reason, we also keep unresolved stores unresolved for 5000 cycles, and give the load and the store the same address.

\subsubsection{Interleaving and State Assignment}\label{sec:conc-interleave}
\reviewtag{[TODO: Applied your requests in this subsection; check it]}
We then add a second round of interleaving: we inject one or two synthetic chains of loads, stores and ALU ops at random positions after the committed prefix, with every chain operand drawn from a shared free pool of unspecified state (Fig.~\ref{fig:interleave}). Because a chain's output can feed an unspecified operand of the backbone, interleaving builds transmitters fed by two independent speculative sources, which no single Alloy instance can express;
\notebr
\auditnote{fact}{fact}{This is a capability the evaluated pipeline does not use. Chain interleaving is off in every evaluation config, and on 300 STT\_6 load-transmitter instances it put an ALU on the transmitter's address path 0 of 1,200 times, because that address is always an Alloy-locked register and chains write only free registers. `Which no single Alloy instance can express' is also wrong: the shape is satisfiable when forced; minimality drops the redundant second source, so no Stage-1a core contains it. SimSpect reproduces the getOldestTaint merge with the OTB merge of fanout v2 (flat test set, ReCon-STT build: 100/100 tests hit, 0/100 without it), which is on in the 2026-11-09 evaluation configs. \newline \textbf{PROPOSED TEXT} (edits the sentence just before this note): interleaving can feed a transmitter from two independent speculative sources, a shape that minimal graphs never contain; \newline \textit{Source:} wiki/design/VOCAB\_TABLES.md (interleave, fanout v2 OTB); wiki/findings/mdsToRead/OTB\_GENERATION\_GAP.md (forcing is SAT); runconfigs/twostage\_fv2\_20261109/*.jsonc}
\notebr
chains off the backbone also add program length and timing variety.
We have two sets of state elements — those determined in Alloy, and those left unspecified. We handle the predetermined state by tying physical registers and memory locations to Alloy register and memory atoms. We then reserve those state locations, and draw every unspecified operand from the shared free pool, the same pool the interleaved chains use
\notebr
\auditnote{fact}{fact}{Only unspecified register operands share the pool. Each unspecified memory operand gets a fresh, never-aliased slot. An unspecified address register of a load or store is not emitted: a constant offset is used instead. Without interleaving, the pool holds no chain registers. \newline \textbf{PROPOSED TEXT} (edits the sentence just before this note): draw every unspecified register operand from the shared free pool, the same pool the interleaved chains use (unspecified memory operands get fresh, unaliased slots) \newline \textit{Source:} STAGE2\_compilation/parsexml.py:1601 (pass3: fresh mem slot), :2146 (pass5: only rf-specified inaddr used), pass\_interleave returns empty pools when disabled}
\notebr
so an unspecified input can read an injected chain's output. All registers employed in memory accesses take the form of offsets from a base pointer to prevent segmentation faults due to random accesses to unallocated memory.
Finally, we convert each program to SSA form, emit LLVM IR with a label at every PC, compile it to x86 assembly, and write the annotation that Stage~3 reads (Fig.~\ref{fig:pipeline}, Stage~2, passes~6--7).
\urgentbr
\urgentnote{\textbf{PROPOSED TEXT (how a test survives compilation; insert after the sentence above).} The compiler must not drop, merge, reorder or re-register any instruction the model specifies. Each instruction is an \texttt{asm sideeffect} block, which LLVM never removes or reorders, and every register the model fixes is pinned to its physical register by the block's register constraint (for a branch, only the condition set-up is inline assembly; LLVM emits the compare and the jump). Memory is a single array whose slots are zero-initialised with volatile stores. We compile with clang at \texttt{-O2}: at \texttt{-O0} the register allocator spills every live value at each basic-block boundary despite the pinning, and the extra loads and stores cause spurious memory-order violations on the speculative path. After linking, the harness reads each PC label's address back from the binary, and the direction of each compiled branch jump, and records them in the annotation, so every directive targets the instruction the model names. \textit{Notes:} two details of the old draft (oldtechnicaldetails.tex:88) are out of date: the code no longer uses \texttt{fake.use}, and the register pool is no longer caller-saved only (it adds r12--r15, which LLVM saves and restores). \textit{Source:} STAGE2\_compilation/compile\_annotate.py:86-105 (the -O2 rationale), :137-222 (PC labels from nm, branches from objdump); STAGE2\_compilation/parsexml.py:1428-1439 (register pool), :2097-2100 (volatile zero-initialisation), :2191 (pinned register write); the instruction-choice note above (branches).}
\notebr

\subsection{Gem5 Execution and Instrumentation}
\reviewtag{[TODO: Applied your requests in this subsection; check it]}
Our tests and generated speculation directives must finally be run on the gem5 simulator. The final phase of testing has a few imperatives to ensure that each hit's gem5 trace matches the modeled violation (violation = the model's graph; hit = what gem5 flags). 

It enforces resolution and speculative program order by using the white-box nature of the simulator to actively set this state. However, it enforces transmitter transmission at a time that matches committed annotations by passively measuring the pipeline, and this is how we detect a violation.
\urgentbr
\urgentnote{\textbf{VOCAB}: ``this is how we detect a violation'' $\to$ ``this is how we detect a hit''.}
\notebr
\urgentbr
\urgentnote{\textbf{PROPOSED TEXT (how we use the white-box simulator; from the old draft, oldtechnicaldetails.tex:42 and :98-101).} White-box access lets the harness set each condition of the modeled execution directly instead of coaxing it. (1) \emph{Speculation}: instead of training the predictors, it overrides them at each annotated PC, forcing a branch to mispredict or to predict correctly, and on the store-bypass build forcing a load to predict no alias with an older store. (2) \emph{Timing}: it holds a branch's resolution, and on the store-bypass build a store's address, for a set number of cycles, so the window is wide and its length is known. (3) \emph{Observation}: it reads the outcome from the simulator's debug traces, which give the tick at which each instruction executes, obtains data, resolves and is squashed, so a hit is read off the trace rather than inferred (see Transmission Requirement). The defense's own logic is never modified. \textit{Placement:} this half of old line 42 belongs here, with the hooks. The other half (realizing the same directives on a target without white-box access by pretraining the predictor) is about generality and goes in the Discussion; there is a note there. \textit{Source:} the gem5-hooks paragraph of this section; gem5-hooks/README.md; wiki/design/RESOLUTION\_HOOKS.md sec.~1.}
\notebr

\subsubsection{Speculation Directives}
The simulator must be made to actually realize the modeled speculative scenario --- the intended branches and dataflow must mispredict, the specified branches must stay unresolved through the transmitter's action, and the commit boundary must stay where is specified \emph{without} altering anything the defense does. If the Alloy model describes an abstract scenario (``a branch mispredicts, opening a window; a load executes speculatively inside it''); a stock CPU will not necessarily realize that exact scenario for a given litmus binary (the predictor may guess right, or the window may be too short for the transmitter to fire, even if the defense might permit this). We bridge this gap with a small, fixed set of hooks, partitioned into \emph{control} hooks (C++ source in the gem5, which realize the scenario), \emph{config glue} (Python, which drives the hooks from the per-test annotation), and \emph{observation} (existing debug flags, which expose the transmitter, as described below). The governing constraint is that we touch only these hooks and never the defense's taint/protection logic, and each control hook is paired with a control run --- the fnc-stall on vs.\ off, the branch-annotation on vs.\ off, and the \texttt{Unsafe}-scheme baseline --- proving the hook did not itself manufacture the result.
\notebr
\reviewnote{\textbf{CHECK}: of the three controls, only the Unsafe baseline is in the evaluation (row E1, run on the 2026-11-09 two-stage sample). The fnc commit stall is 0 in every evaluation config, so its on/off pairing does not apply, and no run turns the branch annotation off. Reword to ``can be paired'', or name only the Unsafe baseline. \textit{Source:} runconfigs/twostage\_fv2\_20261109/*.jsonc (fnc\_commit\_stall\_cycles 0); wiki/state/EVAL\_MANIFEST.md rows E1-E14.}
\notebr

% \paragraph{Directives passed from concretization.}
The concretization phase emits, per test, an annotation (\texttt{.ann.json}) carrying: the BTB target each forced branch will receive (where it will predict to); the resolution time of every speculation primitive (the per-branch resolve delay, and the long stall held on unresolved branches and stores); and the (load, store) PC pairs that should incorrectly alias for store bypass.
The same annotation file both tells gem5 how to reproduce the window and tells the checker what to look for.
\notebr
\reviewnote{\textbf{FIX}: Stage~2 writes no timing, and the pairs do not ``incorrectly alias'': the load is forced to mispredict no alias. \newline \textbf{PROPOSED TEXT}: ``\ldots{} carrying each element's PC, each branch's mispredict mode and BTB target, the transmitter's kind, and for store bypass the (store, load) PC pair whose load is forced to mispredict no alias. Stall lengths come from the run configuration.''}
\notebr

\subsubsection{gem5 hooks}
Three source control hooks realize the modeled scenario, each acting only on control or timing (prediction, redirect latency, commit pacing) and never on the defense's data path. A \emph{BTB-forced target} overrides the predictor at an annotated branch PC to force the modeled mispredicted (taken) path.
\notebr
\reviewnote{\textbf{UPDATE}: on the hooked builds every annotated branch gets a forced prediction in either direction: unresolved branches are forced to mispredict, resolved branches to predict correctly (Table~\ref{tab:instrtable}, rows 1--2).}
\notebr
A \emph{per-branch resolve-stall} delays a branch's resolution to slide the speculative window
over the transmitter. A \emph{commit-stall} holds the first uncommitted instruction --- the speculation boundary --- at the ROB head to widen the window. Rather than fix these stalls, we sweep them so that a reported leak is shown not to hinge on any single window setting:
\notebr
\auditnote{fact}{fact, high}{The code does the reverse. A test is reported as leaking if any grid point triggers (OR over the sweep), so a reported leak can depend on a single setting. The per-point outcomes are saved in the *\_sweep.json sidecars (triggered\_points), so robustness could be reported, but the verdict does not mean that. \newline \textbf{PROPOSED TEXT} (edits the sentence just before this note): We sweep these stalls and report a test as leaking if any setting triggers it; per-setting outcomes are kept, so for each leak we also report how many settings trigger it. \newline \textit{Source:} pipeline.py:330-339 (comment: 'aggregates per-test as OR over grid points'), :500-558 (\_aggregate\_sweep\_results: any\_hit, triggered\_points)}
\notebr
the resolved-branch stall over $s_R \in \{0, 500, 2500\}$ cycles, the unresolved-branch stall $s_U$ (a long $5000$-cycle value masks leaks by pushing the transmitter past \texttt{fnc\_retire}, so runs are launched bare at $s_U = 0$), and the commit-stall over $\{0, 150\}$ cycles.
\urgentbr
\urgentnote{These are the pre-fix settings and contradict Table~\ref{tab:instrtable} and the speculation-directives paragraph. The evaluated configs sweep $s_R \in \{0, 500, 2500\}$ (footprint rows and the AMuLeT STT row: 0 only), hold every unresolved branch for $s_U$ cycles (5000 in the 2026-10-30 fixed configs; the 2026-11-09 two-stage configs draw $s_U$ per branch from $\{2500, 5000, 10000\}$), and set no commit stall. The claim that a 5000-cycle $s_U$ masks leaks predates the resolve-stall fix. \textbf{PROPOSED TEXT} (edits the sentence just before this note): ``\ldots{} the resolved-branch stall over $s_R \in \{0, 500, 2500\}$ cycles; unresolved branches are held for $s_U$ cycles, long enough that they stay unresolved past the transmitter, and the commit stall is off (Table~\ref{tab:instrtable}).'' \textit{Source:} runconfigs/fixed\_20261030/STT\_6\_sttmaster\_fixed\_latest.jsonc:47-57; runconfigs/twostage\_fv2\_20261109/E1\_STT\_6\_fv2.jsonc (sweep).}
\notebr
Finally, for data-flow speculation we instrument store bypass in the false-negative direction --- a load fails to find a store in its store set --- by holding the store unresolved and forcing the aliasing load and store to the same address.
\notebr

Realizing these directives requires a small, fixed set of source sites in every defense build, which is also the checklist for porting the pipeline to a new defense: \texttt{pred/BranchPredictor.py} declares the annotation parameters (forced branch PCs, BTB targets, resolve-stall cycles); \texttt{pred/bpred\_unit.\{cc,hh\}} initializes those maps and overrides the branch target and resolve-stall at predict time; \texttt{o3/\{O3CPU,BaseO3CPU\}.py} declares the fnc-commit-stall parameters; \texttt{o3/cpu.\{hh,cc\}} carry the stall into the CPU; \texttt{o3/commit\{.hh,\_impl.hh\}} add the stall loop, before the squash check, that holds the first uncommitted instruction at the ROB head; \texttt{common/Options.py} exposes the four CLI flags; and \texttt{example/se.py} loads the per-test annotation and wires the targets and stalls into each CPU's branch predictor. Every build also carries two later hook sets. The \emph{resolve hooks} hold a branch's resolution (its execute completion, hand-off to commit and squash broadcast) for the annotated number of cycles, force its prediction, wake the CPU when the hold is released, and print the branch's predict and resolve events. On the store-bypass build, the \emph{store-bypass hooks} hold the annotated store's address generation for the configured number of cycles and force the paired load's memory-dependence prediction to no alias.
\reviewtag{[TODO: Added the resolve and STL hooks (gem5-hooks/README.md; RESOLUTION\_HOOKS.md; STL\_SPECULATION.md)]}
\notebr
\reviewnote{\textbf{FIX}: not every build carries both later hook sets. The resolve hooks are on all ten hooked builds; the store-bypass hooks are only on the store-bypass build, which predates the resolve hooks (rebuild planned), and its scorer is a draft, not a window checker. \newline \textbf{PROPOSED TEXT}: ``Every hooked build also carries the \emph{resolve hooks} \ldots{}; the store-bypass build adds the \emph{store-bypass hooks} \ldots'' \textit{Source:} wiki/state/EVAL\_MANIFEST.md STL section; project log 2026-11-08 21:30 (recon-ruby) and 2026-11-09 13:25 (stl-campaign).}

Detecting transmission requires the pipeline to actually \emph{print} the values the checker reads, so we run gem5 with its existing \texttt{O3PipeView}, \texttt{LSQUnit}, and \texttt{Commit} debug flags --- which emit, respectively, per-instruction issue/complete/squash ticks, the load's execute/packet/spec-read ticks, and the pc$\to$squash-tick map --- and parse that trace. No cache or defense code is instrumented to \emph{see} a leak, only to \emph{realize} the scenario;
\notebr
\auditnote{fact}{fact}{SimSpect does add instrumentation whose only job is observation. It adds print-only BR predict/resolve lines in bpred, IEW and commit on all ten hooked builds (the eight of RESOLUTION\_HOOKS.md plus stt-master and recon-ruby3), and cache, TLB and buffer event prints in the Ruby and classic cache code of the InvisiSpec, CleanupSpec and SpecLFB builds. The resolve hooks also change behaviour: the release now wakes the CPU. Also in this paragraph: O3PipeView prints fetch-to-retire ticks but no squash ticks. Squash ticks come from LSQUnit and Commit. \newline \textbf{PROPOSED TEXT} (edits the sentence just before this note): Observation hooks are print-only: they add trace lines (branch prediction and resolution, and on the footprint builds cache and buffer events), verified to leave the rest of the trace unchanged; no defense logic is modified. \newline \textit{Source:} wiki/design/RESOLUTION\_HOOKS.md sec.1 (BR lines; edit 4 wakeCPU); wiki/design/SIMSPECT\_HOOKS.md header; gem5-hooks/README.md table; STAGE3\_gem5/gem5\_common.py (parse\_pipeview) (parse\_pipeview fields)}
\notebr
where a build does not already print a needed signal, we add a minimal print of the missing tick rather than touch any defense logic.

\subsubsection{Transmission Requirement}
Given a target transmitter defined by our leakage backbone, we must specify how to detect transmission. We do this passively: unlike resolution and speculative program order, which we enforce by actively setting state in the white-box simulator, transmission is detected by \emph{measuring} the pipeline, so that we never perturb the defense's own data path. We define detection on a \textbf{per-transmitter pipeline predicate} which takes the gem5 execution trace (the existing \texttt{O3PipeView}, \texttt{LSQUnit}, and \texttt{Commit} debug streams, parsed by our harness) and analyzes it for the transmitter's characteristic signal. If that signal occurs during the \textbf{violation window}, we report a hit.
Fig.~\ref{fig:checkerwindow} shows this decision on a real gem5 trace.
\urgentbr
\urgentnote{\textbf{PROPOSED TEXT (no attacker or probe needed; from the old draft, oldtechnicaldetails.tex:107).} Because we read the simulator's trace directly, detection needs no attacker: we neither time a probe nor compare two runs. The attacker model is already part of the specification. The leakage contract names the transmitters and operands the defense treats as observable, so a hit is simply one of those transmissions inside the violation window: for a load, the transmitter obtained its data; for the footprint defenses, a squashed instruction changed a structure the defense protects. AMuLeT, by contrast, runs a program on pairs of inputs whose contract traces match and flags a difference in the microarchitectural state it observes~\cite{amulet}. \textit{Source:} oldtechnicaldetails.tex:107; background.tex:23 and the AMuLeT note in intro.tex (fu-2025-amulet.pdf \S2.3); STAGE3\_gem5/check\_ld.py (data-obtained gate); the footprint note on Fig.~\ref{fig:checkerwindow} below.}
\notebr

\notebr
\reviewnote{Pick one window name: the checker figure's caption calls it the unresolved window and ends it at the branch's resolution.}
\notebr

\input{tikz/checkerwindow_v4_8pt}% moved up to where it is first cited, 2026-11-03
\urgentbr
\urgentnote{\textbf{VOCAB}: In Fig.~\ref{fig:checkerwindow} (file tikz/checkerwindow\_v4\_8pt.tex), relabel ``Unresolved window'' to ``Violation window''. In its caption, change ``How the checker decides a violation'' to ``How the checker decides a hit'', ``the unresolved window runs from'' to ``the violation window runs from'', and ``reported as a violation'' to ``reported as a hit''.}
\notebr
\auditnote{fig}{figure, high}{\textbf{About the caption of Fig.~\ref{fig:checkerwindow}:} This example is now in scope. Since 2026-11-08 each footprint criterion checks only the structures its defense's paper claims (for InvisiSpec, MSHR contention and the I-cache are out, the L1D is in), and UV1 is found in 2,995 of 3,000 tests on the hooked InvisiSpec build. What is out of date is the results section's scope paragraph, if it still calls UV1 collateral. \newline \textbf{PROPOSED TEXT}: keep this example, caption it as InvisiSpec's footprint criterion, and update the scope paragraph in the results. \newline \textit{Source:} wiki/findings/mdsToRead/AMULET\_BUGS\_DETECTION\_ACCOUNTING.md sec.~2 (UV1 row); specs/leakage/invisispec.json}
\notebr
\auditnote{claim}{claim}{\textbf{About the caption of Fig.~\ref{fig:checkerwindow}:} check\_footprint counts a footprint by any squashed instruction younger than lc, not just the transmitter. In this trace seqs 566, 573 (xmit) and 576 all qualify, and the earliest in-window footprint is seq 566's L2 coherence change at cycle 16, so the verdict does not depend on xmit's eviction. The fixture also predates the 2026-10-30 harness fix (5000-cycle resolve stall).  \textbf{PROPOSED TEXT} (edits the sentence just before this note): xmit's L1D eviction (\$\textbackslash{}star\$) is one of the in-window footprints that meet InvisiSpec's criterion (the criterion counts any squashed load younger than lc), so the test is reported as a violation.  \textit{Source:} project log 2026-11-02 'InvisiSpec footprint checker attributes leaks to any squashed load'; paper\_figures/plans/fig\_e\_checkerwindow.md:140,199; URGENT\_TODO.md InvisiSpec footprint DECIDE}
\notebr

The violation window
\notebr
is the interval during which the \emph{last committed} instruction (\texttt{lc}) has committed but the \emph{first uncommitted} instruction (\texttt{fnc}) has not. Because commit is in order, this is the interval \texttt{(lc\_retire, fnc\_retire)}, with the additional requirement that all branches guarding the transmitter are still unresolved at the transmission tick.
\notebr
\reviewnote{\textbf{REWORD} ``all branches guarding the transmitter'' to ``the branches the model marks unresolved'': that is the gate the checker applies.}
\notebr
A transmission inside this window, on the modeled mis-speculated path, is a violation;
\urgentbr
\urgentnote{\textbf{VOCAB}: ``is a violation;'' $\to$ ``is a hit;'' (a gem5 result; it becomes a leak only once validated).}
\notebr
the same action after the window closes is not (Figure~\ref{fig:checkerwindow}).
We call this the \emph{violation window} because it is where every $\mu$arch condition the instance specifies still holds (which instructions have committed and which branches are unresolved), not only the commit status of \texttt{lc} and \texttt{fnc}; a transmission inside it is exactly what the leakage predicate forbids.

% \input{tikz/xmitcriteria_8pt}% hidden 2026-11-03 (Anna); file kept
% \figtodo{transmission-criteria table, Table~\ref{tab:xmitcriteria}: reference it in the transmission-requirement text; its content is still under review}

\urgentbr
\urgentnote{\textbf{VOCAB}: \textbf{ADD DEFINITION} (new text, insert here): ``We use three terms for three stages. A \emph{violation} is an abstract execution graph that breaks the defense's leakage predicate (Sec.~III). A \emph{hit} is a test whose transmitter the checker sees transmit inside the violation window on the gem5 build. A \emph{leak} is a hit we confirm is real, and a \emph{bug} when the defense should have blocked it.''}
\notebr
What counts as ``a transmission'' is transmitter-specific, and crucially \emph{different transmitter types transmit through different pipeline events, at different pipeline stages, and therefore at different times}. We thus need a distinct signal per transmitter:
\begin{itemize}
  \item \textbf{Load transmitter.} The signal is the load's cache touch --- the tick at which its access reaches the cache
\notebr
\auditnote{fact}{fact}{check\_ld times the load by its latest LSQ `Executing load' tick, which fires before the defense decides whether to send the access; the pipeview complete tick is the fallback. The packet and spec\_read ticks only feed the yes/no data-obtained gate and are never compared with the window. If a multi-micro-op transmitter (a read-modify-write) never executes its load micro-op, the fallback scores the first micro-op's complete tick. The footprint defenses (InvisiSpec, CleanupSpec, SpecLFB) do not use this predicate: their runs use the footprint check over SimSpect cache-event lines with a per-defense criterion, and the memaccess and expose checkers are legacy. \newline \textbf{PROPOSED TEXT} (edits the sentence just before this note): The signal is the tick at which the LSQ executes the load (\texttt{Executing load}; \texttt{O3PipeView} \texttt{complete} if absent), and the load counts only if it then obtained data \ldots; the footprint defenses use their own footprint criteria. \newline \textit{Source:} STAGE3\_gem5/check\_ld.py:10-21, 83-87, 119-122; runconfigs/footprint/run\_config\_STT\_6\_speclfb\_footprint.jsonc:44 (check\_footprint), :68 (legacy\_checker)}
\notebr
--- read from the \texttt{LSQUnit} trace (the \texttt{Executing load} / packet / \texttt{spec\_read} ticks; on newer builds an additive \texttt{memaccess} tick for SpecLFB and an \texttt{expose} tick for InvisiSpec) together with the \texttt{O3PipeView} \texttt{complete} tick. A load counts as transmitting only if it is in-window \emph{and} it actually \textbf{obtained data} --- either it \emph{reached the cache} or it received a \emph{store-to-load forward}. Squash timing never gates a load: one that obtained data in the window has already transmitted even if it is later squashed.
   % This data-obtained gate is essential: store-to-load forwarding lets a speculative load complete and leak \emph{without} a cache packet, so a cache-reach-only signal would miss it.

  

  \item \textbf{Branch transmitter.} The signal is the branch's redirect, keyed on the \emph{squash} tick (the tick at which the misprediction broadcasts a redirect), read from the \texttt{Commit} trace as a pc~$\rightarrow$~squash-tick map. A branch transmits later in the pipeline than a load (at resolution/redirect rather than at execute), which is precisely why it needs its own predicate and its own timing.

  \item \textbf{Faulting / long-latency ``other'' transmitter.} A speculative transmitter that faults (e.g.\ a divide) transmits through its timing: the signal is its completion tick, and a fault after a mispredicted branch is scored as a transmission (Table~\ref{tab:instrtable}, row~5).
\reviewtag{[TODO: Finished this sentence]}
  % ; a speculative divide fault is promoted to a leak, which is what makes a guard such as the \texttt{idivq} guard meaningful.

  \item \textbf{Store transmitter.} A store transmits through its address-dependent effects --- its TLB/cache access and, for the store-to-load-forward channel, the forward it provides to a younger load. A speculative store does not install a cache line until commit, so its in-window signal is the address-dependent access/forward rather than a line install.
\notebr
\reviewnote{There is no store checker: kinds without their own checker fall back to the load check (pipeline.py:313-320), which scores a store by its completion tick. Only our SPT model makes stores transmit (SPT\_6\_oneLP\_v2.als, leakage function), and store transmitters are out of scope for STT (decided 2026-11-08: STT assumes stores act only at retirement). CUT this bullet, or \newline \textbf{PROPOSED TEXT}: ``Store transmitters, used only by our SPT model, are scored by the load predicate's fallback, the store's completion tick.''}
\notebr
\end{itemize}
Because every reported leak must be a transmission inside the window on the modeled path, the predicate is paired with the window gate and the ''branches still unresolved'' gate.
\notebr
\reviewnote{typo: the opening quote should be two backticks.}
\notebr

% , and (for a \emph{demonstrable} leak rather than a dataflow signal) optionally with a data-taint gate requiring the touched address or victim to depend on the tainted operand.

\subsection{The Leakage-prompting $\mu$arch Knobs}
The dimensions the pipeline exposes fall into three classes. Enumeration is \emph{comprehensive} over the \emph{forms of leakage} --- the minimal, necessary leakage backbones the defense recognizes --- so coverage is complete up to the instruction bound. The concretization choices that merely \emph{realize} a backbone --- x86 instruction templates, the two branch-misprediction modes, and the interleaved chains --- are \emph{semi-enumerated}, sampled rather than exhausted. The leakage predicate is fixed: the Alloy predicates that define a violation, and, per campaign, the defense under test. 
\notebr
% --- is fixed, while every knob that \emph{probes} it, from the window stalls to the micro-architectural timing of the fragility sweep, is reported as a swept variable, never tuned to hit or avoid a bug.
\anna{idk}